\documentclass[11pt,journal,onecolumn]{IEEEtran}

\usepackage{geometry}
\geometry{letterpaper, margin=1in}
\usepackage{graphicx}
\usepackage{booktabs}
\usepackage{hyperref}
\usepackage{amsmath}
\usepackage{microtype}
\usepackage{xcolor}
\usepackage{subcaption}

\hypersetup{
    colorlinks=true,
    linkcolor=blue,
    filecolor=magenta,      
    urlcolor=cyan,
    pdftitle={Multimodal Emotion Recognition on TESS},
    pdfpagemode=FullScreen,
}

\title{Multimodal Emotion Recognition on the TESS Dataset}
\author{Aaron Roy}
\date{\today}

\begin{document}

\maketitle

\begin{abstract}
This report outlines the architecture decisions, experiments, and results of a Multimodal Emotion Recognition system built on the Toronto emotional speech set (TESS) dataset. We evaluate three pipeline variants: Speech-only (Mel Spectrograms + BiLSTM), Text-only (Pre-trained BERT Embeddings), and Multimodal Fusion (Late Fusion via concatenation). The fusion model achieves the highest overall classification accuracy of 99.38\%, highlighting the synergy between acoustic feature modeling and contextual text grounding.
\end{abstract}

\section{Introduction}
Emotion recognition from human conversation is a crucial challenge in human-computer interaction (HCI). While voice tone (speech prosody) carries significant acoustic cues, the literal words spoken (linguistic text context) provide critical constraints that resolve phonetic and acoustic ambiguities. In this project, we develop and compare three deep learning architectures: a Speech-only model, a Text-only model, and a Multimodal Fusion model combining both modalities.

\section{Architecture Decisions}
The overall system is designed around five modular blocks: preprocessing, feature extraction, temporal/contextual modeling, feature fusion, and classification.

\subsection{Preprocessing}
\begin{itemize}
    \item \textbf{Speech Modality:} We use \texttt{torchaudio} to load raw audio files and ensure a standard 16 kHz sampling rate via resampling. Waveforms are padded or truncated to exactly 64,000 samples (4.0 seconds) to maintain consistent batch shapes.
    \item \textbf{Text Modality:} Since the TESS dataset carrier phrases follow a strict folder structure and file naming convention (e.g., ``Say the word [word]''), we derive transcripts directly from the filenames. We tokenize these carrier phrases using the pre-trained \texttt{bert-base-uncased} tokenizer with a maximum length of 16 tokens.
\end{itemize}

\subsection{Feature Extraction}
\begin{itemize}
    \item \textbf{Speech Features:} We extract Mel Spectrograms (64 mel bins, FFT size of 1024, hop length of 512) from the resampled audio waveform. Mel Spectrograms represent spectral energy on a logarithmic scale mimicking human auditory perception.
    \item \textbf{Text Features:} We pass tokenized input IDs and attention masks through a pre-trained BERT model to extract dense contextual word embeddings.
\end{itemize}

\subsection{Temporal / Contextual Modeling}
\begin{itemize}
    \item \textbf{Speech Modality:} We employ a \textbf{Bidirectional Long Short-Term Memory (BiLSTM)} network with 2 layers and 128 hidden units. The bidirectional structure allows the network to process sequential acoustic features from both past and future temporal context. The final hidden states of the forward and backward LSTMs are concatenated to form a 256-dimensional summary vector.
    \item \textbf{Text Modality:} We extract the pooled output (corresponding to the \texttt{[CLS]} token) from the pre-trained \texttt{bert-base-uncased} model. The BERT weights are frozen during training to ensure stability and efficiency, producing a 768-dimensional contextual document embedding.
\end{itemize}

\subsection{Fusion Block}
We implement a \textbf{Late Fusion} approach. The 256-dimensional representation from the Speech BiLSTM and the 768-dimensional embedding from BERT are concatenated directly to form a unified 1024-dimensional multimodal representation. Late fusion allows the speech and text pipelines to specialize in feature extraction independently before mapping combined high-level concepts to emotion labels.

\subsection{Classifier}
Each pipeline variant (Speech, Text, and Fusion) terminates with a classifier head consisting of fully connected linear layers, Dropout regularization (rate = 0.3) to prevent overfitting, and a Softmax output layer representing the 7 emotion classes.

\section{Experiments and Performance}
We trained the models on an 80/20 train/test split of the TESS dataset across 10 epochs with a batch size of 32 using the Adam optimizer (learning rate = 0.001). The dataset contains 2,800 audio samples distributed across 7 emotion categories: \textit{angry, disgust, fear, happy, pleasant surprise, sad, and neutral}.

Table~\ref{tab:overall_perf} outlines the overall test accuracy across the three pipeline variants. Table~\ref{tab:per_emotion_f1} lists the individual F1-scores per emotion.

\begin{table}[h]
\centering
\caption{Overall Model Performance Comparison}
\label{tab:overall_perf}
\begin{tabular}{lll}
\toprule
\textbf{Model Variant} & \textbf{Modalities Used} & \textbf{Overall Test Accuracy} \\
\midrule
Speech-Only & Mel-Spectrograms & 97.32\% \\
Text-Only & BERT Embeddings & 13.48\% \\
\textbf{Multimodal (Fusion)} & \textbf{Speech + Text} & \textbf{99.38\%} \\
\bottomrule
\end{tabular}
\end{table}

\begin{table}[h]
\centering
\caption{Per-Emotion Classification F1-Score Breakdown}
\label{tab:per_emotion_f1}
\begin{tabular}{lccc}
\toprule
\textbf{Emotion} & \textbf{Speech-Only F1} & \textbf{Text-Only F1} & \textbf{Multimodal Fusion F1} \\
\midrule
Angry & 1.00 & 0.00 & 1.00 \\
Disgust & 0.93 & 0.00 & 0.99 \\
Fear & 1.00 & 0.00 & 1.00 \\
Happy & 0.98 & 0.08 & 0.99 \\
Pleasant Surprise & 0.98 & 0.23 & 0.98 \\
Sad & 0.94 & 0.00 & 0.99 \\
Neutral & 0.99 & 0.10 & 1.00 \\
\bottomrule
\end{tabular}
\end{table}

\section{Analysis}

\subsection{Emotion Classification Difficulty}
In the Speech-only model, \textit{fear}, \textit{angry}, and \textit{neutral} were the easiest emotions to classify, yielding F1-scores between 0.99 and 1.00. This is attributed to their distinct acoustic properties (e.g., elevated pitch and intensity for anger/fear; flat pitch and speech rate for neutral). 

Conversely, the hardest emotions to classify solely by audio were \textit{disgust} (F1 = 0.93) and \textit{sad} (F1 = 0.94). These emotions possess low arousal profiles, lower voice intensity, and slower speech rates, causing spectral representation overlaps that make acoustic separation difficult.

\subsection{The Value of Modality Fusion}
The Text-only model performs at 13.48\%, which is equivalent to random guessing over 7 classes. This is mathematically expected: because the TESS transcripts are identical carrier phrases (``Say the word [word]''), the text contains absolutely no semantic emotional cues. 

However, late fusion improves overall accuracy from 97.32\% (Speech-only) to \textbf{99.38\%}. While the text itself contains no emotional valence, the BERT embedding informs the model of the \textit{exact identity} of the target word spoken. The fusion network learns to evaluate acoustic variations conditioned on the specific phonetic structure of the word being spoken, eliminating acoustic ambiguities.

\subsection{Error Analysis and Failure Cases}
\begin{itemize}
    \item \textbf{Text-Only Failures:} The model suffers from complete classification collapse, predicting \textit{pleasant surprise} or \textit{neutral} for almost all test samples due to a lack of emotional features in the carrier phrase transcripts.
    \item \textbf{Speech-Only Confusions:} The majority of confusion occurs between \textit{disgust} and \textit{sadness}. Because both are low-intensity, negative-valence emotions, their acoustic Mel-spectrogram profiles are highly similar. Late fusion resolves this confusion by introducing targets vocabulary context.
\end{itemize}

\subsection{Cluster Separability Visualization}
We generated t-SNE scatter plots mapping the high-dimensional feature vectors into a 2D space to visualize the learned representations across different stages of the model:
\begin{itemize}
    \item \textbf{Speech Temporal Modeling} (Figure~\ref{fig:speech_tsne}): Features extracted from the BiLSTM network demonstrate clear spatial separation, indicating successful representation learning.
    \item \textbf{Text Contextual Modeling} (Figure~\ref{fig:text_tsne}): Features extracted from BERT form a single overlapping cluster, confirming a lack of intrinsic emotional differentiators.
    \item \textbf{Multimodal Fusion} (Figure~\ref{fig:fusion_tsne}): Features from the fusion layer show extremely tight, separated, and distinct clusters for all 7 emotions, correlating with the near-100\% classification accuracy.
\end{itemize}

\begin{figure}[htbp]
    \centering
    \begin{subfigure}[b]{0.32\textwidth}
        \centering
        \includegraphics[width=\textwidth]{speech_tsne.png}
        \caption{Speech t-SNE}
        \label{fig:speech_tsne}
    \end{subfigure}
    \hfill
    \begin{subfigure}[b]{0.32\textwidth}
        \centering
        \includegraphics[width=\textwidth]{text_tsne.png}
        \caption{Text t-SNE}
        \label{fig:text_tsne}
    \end{subfigure}
    \hfill
    \begin{subfigure}[b]{0.32\textwidth}
        \centering
        \includegraphics[width=\textwidth]{fusion_tsne.png}
        \caption{Fusion t-SNE}
        \label{fig:fusion_tsne}
    \end{subfigure}
    \caption{t-SNE Separability Visualizations for each Modality}
    \label{fig:tsne_plots}
\end{figure}

\begin{figure}[htbp]
    \centering
    \begin{subfigure}[b]{0.32\textwidth}
        \centering
        \includegraphics[width=\textwidth]{speech_confusion_matrix.png}
        \caption{Speech Confusion Matrix}
        \label{fig:speech_cm}
    \end{subfigure}
    \hfill
    \begin{subfigure}[b]{0.32\textwidth}
        \centering
        \includegraphics[width=\textwidth]{text_confusion_matrix.png}
        \caption{Text Confusion Matrix}
        \label{fig:text_cm}
    \end{subfigure}
    \hfill
    \begin{subfigure}[b]{0.32\textwidth}
        \centering
        \includegraphics[width=\textwidth]{fusion_confusion_matrix.png}
        \caption{Fusion Confusion Matrix}
        \label{fig:fusion_cm}
    \end{subfigure}
    \caption{Confusion Matrices for individual Modalities and Fusion model}
    \label{fig:confusion_matrices}
\end{figure}

\end{document}
